\section{Implementation}
\label{ch:implementation}

This section shows how the design from Section~\ref{sec:design} appears in the
codebase. It connects the class, sequence, component, and data design
artifacts to the implemented backend and frontend modules so that the
design decisions can be checked against the actual source code.

\subsection{Development environment and technology stack}

The project is implemented as a monorepo with separate backend and
frontend runtimes. The backend uses Python~3 and FastAPI~0.133.1
\cite{python2026setup,fastapi2026docs}, SQLAlchemy~2.0.48 and
Alembic~1.7.7 \cite{sqlalchemy2026docs,alembic2026docs},
ChromaDB~1.5.5, OpenAI~2.26.0, and Google Cloud Vision~3.11.0
\cite{chromadb2026docs,openai2026platform,google2026vision}, and
Pillow~12.0.0, Pydantic~2.12.5, Pydantic Settings~2.12.5, and
Uvicorn~0.41.0
\cite{pillow2026docs,pydantic2026docs,pydanticsettings2026docs,uvicorn2026docs}.
The frontend uses React~19.2.0, React Router DOM~7.13.1, and
Vite~7.3.1 \cite{react2026docs,reactrouter2026docs,vite2026docs}.
These exact versions are defined in \texttt{backend/requirements.txt}
and \texttt{frontend/package.json}, which makes the implementation
reproducible instead of describing only a generic stack.

The development environment was intentionally lightweight and local.
Source code was edited primarily in Visual Studio Code
\cite{vscode2026docs}, while Git \cite{git2026docs} was used for
version control across both frontend and backend changes. The backend
runtime was managed through a Python virtual environment
and the dependency manifest in
\texttt{backend/requirements.txt}. The frontend runtime was managed
through Node.js and npm \cite{node2026docs,npm2026docs} and the
dependency manifests \texttt{frontend/package.json} and
\texttt{frontend/package-lock.json}. Database schema evolution was
handled through Alembic migration files.

The technology choices follow the layered architecture from
Section~3.2. FastAPI supports thin route modules with built-in request
validation. SQLAlchemy and Alembic realize the relational model in
code. React and Vite support the page-oriented frontend structure, and
ChromaDB makes retrieval part of the working system instead of a design
idea described only on paper.

\subsection{Repository structure and code organization}

The repository structure makes the architectural decomposition visible
before reading individual modules:

\begin{lstlisting}[caption={Repository structure of the implemented system},label={lst:repo-tree}]
ai-cookbook-thesis/
|-- backend/
|   |-- main.py
|   |-- config.py
|   |-- database.py
|   |-- dependencies.py
|   |-- models.py
|   |-- schemas.py
|   |-- security.py
|   |-- rate_limit.py
|   |-- routes/
|   |   |-- auth.py
|   |   |-- ocr.py
|   |   |-- rag_debug.py
|   |   `-- recipes.py
|   |-- services/
|   |   |-- embedding_service.py
|   |   |-- llm_service.py
|   |   |-- ocr_image_processing.py
|   |   |-- ocr_provider.py
|   |   |-- ocr_service.py
|   |   |-- recipe_image_provider.py
|   |   |-- recipe_image_service.py
|   |   |-- rag_service.py
|   |   `-- vector_store_service.py
|   |-- scripts/
|   |   |-- prepare_recipes.py
|   |   `-- index_recipes.py
|   |-- tests/
|   `-- requirements.txt
|-- frontend/
|   |-- package.json
|   `-- src/
|       |-- App.jsx
|       |-- context/AuthContext.jsx
|       |-- lib/api.js
|       |-- components/NutritionGrid.jsx
|       `-- pages/
|           |-- GenerateRecipePage.jsx
|           |-- NutritionPage.jsx
|           |-- SavedRecipesPage.jsx
|           |-- RecipeHistoryPage.jsx
|           `-- LoginPage.jsx
`-- alembic/
    |-- env.py
    `-- versions/
\end{lstlisting}

Table~\ref{tab:repo-directories} summarizes the purpose of the main
directories and shows how the repository structure reflects the
implemented architecture.

\begin{table}[H]
	\centering
	\scriptsize
	\renewcommand{\arraystretch}{1.2}
	\begin{tabular}{|p{3.4cm}|p{10.5cm}|}
		\hline
		\textbf{Directory or area} & \textbf{Purpose in the project} \\
		\hline
		\texttt{backend/} & Contains the server-side application, including routing, validation, business logic, persistence, security, and integration with retrieval and external AI services. \\
		\hline
		\texttt{backend/routes/} & Defines the HTTP entry points for authentication, recipe generation, recipe persistence, OCR extraction, and debugging-oriented retrieval access. \\
		\hline
		\texttt{backend/services/} & Holds the multi-step workflow logic for recipe generation, retrieval, OCR, image handling, and provider abstraction so that route handlers remain thin. \\
		\hline
		\texttt{backend/scripts/} & Contains offline preparation utilities used to clean the recipe corpus and build the persistent retrieval index. \\
		\hline
		\texttt{backend/tests/} & Stores automated backend tests for routes, validation, security behavior, and service-level logic. \\
		\hline
		\texttt{frontend/} & Contains the browser-facing application, its dependency manifest, and the React source used for the main user workflows. \\
		\hline
		\texttt{frontend/src/pages/} & Implements the top-level task pages such as recipe generation, OCR extraction, saved recipes, history, and login/account management. \\
		\hline
		\texttt{frontend/src/components/} & Stores reusable presentation elements shared across pages, such as structured result display components. \\
		\hline
		\texttt{frontend/src/context/} & Maintains shared frontend state, most importantly authentication state and current-profile information. \\
		\hline
		\texttt{frontend/src/lib/} & Groups client-side utility code for API access and URL construction so that request logic is not duplicated across pages. \\
		\hline
		\texttt{alembic/} and \texttt{alembic/versions/} & Store migration configuration and revision history for relational schema changes, keeping database evolution explicit and reproducible. \\
		\hline
	\end{tabular}
	\caption{Purpose of the main repository directories}
	\label{tab:repo-directories}
\end{table}

This file system mirrors the layered design from Section~\ref{sec:design}. Route
modules act as a thin API boundary, service modules contain business
logic and provider orchestration, schema classes define validated data
contracts, and ORM models implement persistence. This is also the point
where the service-orchestration boundary described in the design chapter
becomes visible:
\texttt{routes/recipes.py} exposes one endpoint to the client, while
multiple service modules collaborate behind that boundary.

\subsubsection{Development workflow and repository discipline}

The development workflow was organized around small, feature-oriented
changes rather than large monolithic rewrites. Backend services, API
routes, migrations, frontend pages, and tests were versioned together in
the same repository so that every non-trivial feature could be traced as
one coherent change set. In practice, this was especially important for
features that crossed multiple layers, such as profile-default-aware
recipe generation or OCR history persistence, because they required
coordinated updates in route logic, schemas, frontend state handling,
and tests.

Version-control work followed the same feature-oriented structure.
Changes were developed as focused units around one task or one related
set of tasks instead of mixing unrelated edits into the same revision.
For example, a persistence-related change typically combined the route
logic, schema updates, frontend integration, and migration file needed
for that feature. This kept the history easier to review and reduced the
risk of partial changes that would compile but leave one layer out of
sync with the others.

Dependency and environment management were also repository-driven. The
backend dependency set is declared in
\texttt{backend/requirements.txt}, frontend dependencies in
\texttt{frontend/package.json}, and environment-dependent settings are
centralized through the backend configuration layer and the
\texttt{.env} file format. Database evolution is versioned explicitly
through Alembic revision files instead of implicit schema changes. This
workflow keeps the implementation reproducible: a given repository state
contains not only application code, but also the dependency manifests,
migration history, and indexing scripts needed to reconstruct the
working system.

\subsection{Core application logic and complex features}

This section focuses on the implementation of the main complex features
promised in the project proposal. The first two are the structured
recipe-generation workflow and the OCR-based nutrition-label workflow.
The third, data persistence and backend integration, is covered in
detail later in this section because it spans relational schema design,
ORM mapping, migration history, and query behavior.

\subsubsection{Application bootstrap and request entry points}

The backend entry point is \texttt{backend/main.py}. It does more than
create a FastAPI object: it prepares the runtime environment, installs
cross-cutting middleware, mounts media storage, and registers feature
routers \cite{fastapi2026docs}. This directly realizes the component
diagram from Figure~3.2.

\begin{lstlisting}[language=Python,caption={FastAPI bootstrap in \texttt{backend/main.py}},label={lst:main-bootstrap}]
app = FastAPI(title="AI Cookbook API")
Path("backend/media").mkdir(parents=True, exist_ok=True)
Path(settings.recipe_thumbnail_directory).mkdir(parents=True, exist_ok=True)
Path(settings.ocr_upload_directory).mkdir(parents=True, exist_ok=True)

app.add_middleware(
    CORSMiddleware,
    allow_origins=[settings.frontend_origin],
    allow_credentials=True,
    allow_methods=["*"],
    allow_headers=["*"],
)
app.add_middleware(
    RateLimitMiddleware,
    default_rule=RateLimitRule(
        max_requests=settings.default_rate_limit_requests,
        window_seconds=settings.default_rate_limit_window_seconds,
    ),
    route_rules={
        "/signup": RateLimitRule(...),
        "/generate-recipe": RateLimitRule(...),
        "/ocr/extract": RateLimitRule(...),
    },
)

app.mount("/media", StaticFiles(directory="backend/media"), name="media")
app.include_router(auth.router)
app.include_router(recipes.router)
app.include_router(ocr.router)
app.include_router(rag_debug.router)
\end{lstlisting}

This bootstrap code mainly provides the shared runtime context for the
feature-specific workflows discussed next. It registers the route groups
and installs the middleware and media handling that those workflows rely
on.

\subsubsection{User-specific defaults and authenticated request context}

Authentication is implemented in three modules that work together:
\texttt{backend/security.py} contains Argon2-backed password hashing
and JWT helpers \cite{argon22026docs,pythonjose2026docs,jwt2015rfc7519},
\texttt{backend/dependencies.py} resolves the current user from the
token, and \texttt{backend/routes/auth.py} exposes the HTTP endpoints.
The hashing and verification calls are implemented as follows:

\begin{lstlisting}[language=Python,caption={Password hashing utilities in \texttt{backend/security.py}},label={lst:auth-hashing}]
password_hash = PasswordHash.recommended()

def hash_password(password: str) -> str:
    return password_hash.hash(password)

def verify_password(plain_password: str, hashed_password: str) -> bool:
    return password_hash.verify(plain_password, hashed_password)
\end{lstlisting}

The dependency that enforces ownership of user-specific data is defined
in \texttt{backend/dependencies.py}:

\begin{lstlisting}[language=Python,caption={JWT-based current-user dependency},label={lst:current-user-dependency}]
oauth2_scheme = OAuth2PasswordBearer(tokenUrl="/token")

def get_current_user(
    token: str = Depends(oauth2_scheme),
    db: Session = Depends(get_db),
) -> User:
    user = _get_user_from_token(token, db)
    if user is None:
        raise credentials_exception
    return user
\end{lstlisting}

This function is the implementation-level counterpart of the ownership
relations in the use-case and class diagrams. Endpoints for saving
recipes, browsing history, and saving OCR results depend on the
resolved \texttt{User} instance instead of trusting client-supplied
identifiers.

Profile defaults are persisted in \texttt{backend/routes/auth.py} and
later reused by the frontend when generation fields are left empty:

\begin{lstlisting}[language=Python,caption={Profile update persistence in \texttt{backend/routes/auth.py}},label={lst:profile-update}]
@router.put("/me", response_model=UserResponse)
def update_current_user_profile(...):
    current_user.preferences = json.dumps(profile.preferences)
    current_user.allergies = json.dumps(profile.allergies)
    db.add(current_user)
    db.commit()
    db.refresh(current_user)
    return _user_to_response(current_user)
\end{lstlisting}

\subsubsection{Structured recipe generation and personalization}

Recipe generation is the most complex workflow in the system. It
combines request validation, safety filtering, retrieval, prompt
construction, output verification, optional image generation, and
persistence. The route in \texttt{backend/routes/recipes.py} acts as a
thin facade: it delegates generation to the language-model service and
thumbnail generation to the image service, then optionally persists the
result for the authenticated user.

Personalization enters the workflow at more than one point. User
preferences and allergies can arrive directly in the request or through
stored profile defaults, and they influence both the initial filtering
step and the later prompt construction.

The first non-trivial step is conflict detection. The service checks
whether user-provided ingredients violate allergy or preference
constraints before retrieval or generation starts:

\begin{lstlisting}[language=Python,caption={Conflict analysis in \texttt{backend/services/llm\_service.py}},label={lst:conflict-analysis}]
allergy_conflicts = find_conflicting_ingredients(
    request.ingredients,
    request.allergies,
)
preference_conflicts = find_preference_conflicting_ingredients(
    request.ingredients,
    request.preferences,
)
safe_ingredients = _remove_conflicting_ingredients(
    request.ingredients,
    conflicting_ingredients,
)

if not safe_ingredients:
    raise HTTPException(
        status_code=status.HTTP_400_BAD_REQUEST,
        detail="All provided ingredients conflict with the selected allergies or dietary preferences.",
    )
\end{lstlisting}

This safety-oriented preprocessing is the concrete implementation of the
restriction logic discussed in Section~\ref{sec:design}. It is also one of the points
where the sequence diagram can be verified against real code.

After filtering, the backend retrieves grounded context from the Chroma
collection \cite{chromadb2026docs}. The retrieval call is implemented
in \texttt{backend/services/rag\_service.py}, and the query vector is
created through the embeddings service \cite{openai2026embeddings}:

\begin{lstlisting}[language=Python,caption={Vector retrieval query in \texttt{backend/services/rag\_service.py}},label={lst:rag-query}]
collection = get_collection()
query_text = build_retrieval_query(request)
query_embedding = create_embeddings([query_text])[0]
response = collection.query(
    query_embeddings=[query_embedding],
    n_results=settings.rag_candidate_count,
    include=["metadatas", "distances"],
)
\end{lstlisting}

The prompt sent to the LLM through the OpenAI API
\cite{openai2026api} is not generic. It is assembled from the
filtered ingredient list, user constraints, and retrieved recipe
context. A shortened version of the prompt builder shows this clearly:

\begin{lstlisting}[language=Python,caption={Prompt construction in \texttt{backend/services/llm\_service.py}},label={lst:recipe-prompt}]
return f"""
You are generating a new recipe for an AI cookbook application.

User ingredients: {ingredient_list}.
Dietary preferences: {", ".join(request.preferences) if request.preferences else "none"}.
Allergies to avoid: {", ".join(request.allergies) if request.allergies else "none"}.

Retrieved recipe context (use only as structural inspiration, do NOT copy ingredients):
{context_block}

Return only valid JSON with this exact structure:
{{
  "title": "string",
  "ingredients": ["string", "string"],
  "preferences": ["string"],
  "allergies": ["string"],
  "steps": ["string", "string", "string"],
  "nutrition_estimate": {{
    "calories": 0,
    "protein": "string",
    "carbs": "string",
    "fat": "string"
  }}
}}
""".strip()
\end{lstlisting}

The implementation does not trust the first model response blindly. It
checks for restriction violations and performs one correction cycle when
necessary:

\begin{lstlisting}[language=Python,caption={Output validation and retry loop in \texttt{backend/services/llm\_service.py}},label={lst:retry-loop}]
messages: list[dict] = [{"role": "user", "content": prompt}]

for _ in range(2):
    response = client.chat.completions.create(
        model=settings.openai_model,
        messages=messages,
        temperature=0.7,
    )
    recipe_text = _strip_code_fences(response.choices[0].message.content or "")
    recipe_json = json.loads(recipe_text)
    last_violations = _find_output_restriction_violations(recipe_json, restrictions)

    if not last_violations:
        additional = _find_ingredient_drift(recipe_json, safe_ingredients)
        recipe_json["additional_ingredients"] = additional
        recipe_json["warnings"] = warnings
        return RecipeResponse.model_validate(recipe_json)

    violation_note = (
        "The previous output violated these rules and must be corrected: "
        + "; ".join(last_violations)
    )
    messages.append({"role": "assistant", "content": recipe_text})
    messages.append({"role": "user", "content": violation_note})
\end{lstlisting}

This retry loop is the strongest implementation evidence that recipe
generation is realized as a fixed multi-stage workflow rather than a
single opaque model call. Validation, retrieval, prompt construction,
output checking, corrective retry, and response shaping appear as
explicit stages in code, which is the implementation-level counterpart
of the workflow diagrams and design rationale from Section~\ref{sec:design}.

The implementation also includes dedicated RAG inspection endpoints in
\texttt{backend/routes/rag\_debug.py}. These endpoints are not part of
the normal browser workflow, but they were useful during development
because they expose the same retrieval pipeline that recipe generation
uses internally. The \texttt{/rag/debug} endpoint returns the generated
retrieval query text, the candidates returned by Chroma, the
allergy-conflict filtering result, the ranking score assigned to each
candidate, and which items were ultimately selected as context. The
\texttt{/rag/collection-stats} endpoint provides a lighter diagnostic
view over the indexed corpus. In practice, these endpoints made it
possible to verify that failures or weak recipe outputs came from
retrieval quality, filtering, or prompt behavior rather than from a
single opaque model response.

\subsubsection{OCR-based nutrition label processing}

The OCR subsystem is implemented as a two-stage pipeline: optical text
recognition followed by semantic structuring into validated nutrition
fields. The abstraction boundary for OCR providers is defined in
\texttt{backend/services/ocr\_provider.py}
\cite{google2026vision}:

\begin{lstlisting}[language=Python,caption={OCR provider strategy in \texttt{backend/services/ocr\_provider.py}},label={lst:ocr-provider}]
class OcrProvider(ABC):
    @abstractmethod
    def extract_text(self, image_bytes: bytes) -> str:
        """Extract raw OCR text from an image payload."""

class GoogleVisionOcrProvider(OcrProvider):
    def extract_text(self, image_bytes: bytes) -> str:
        return extract_text_with_google_vision(image_bytes)

@lru_cache(maxsize=1)
def get_ocr_provider() -> OcrProvider:
    if settings.ocr_provider == "google_vision":
        return GoogleVisionOcrProvider()
\end{lstlisting}

This implements the provider-substitution mechanism described in the
design chapter. The OCR workflow depends on an application-level
interface and can switch concrete providers without changing the route
or higher-level service logic.

Before OCR is executed, images are validated and normalized in
\texttt{backend/services/ocr\_image\_processing.py} with Pillow
\cite{pillow2026docs}:

\begin{lstlisting}[language=Python,caption={Image preprocessing for OCR},label={lst:ocr-preprocess}]
def prepare_image_for_ocr(file_bytes: bytes) -> bytes:
    image = Image.open(BytesIO(file_bytes))
    image.load()

    normalized = ImageOps.exif_transpose(image)
    grayscale = ImageOps.grayscale(normalized)
    contrast_ready = ImageOps.autocontrast(grayscale)
    enlarged = contrast_ready.resize(
        (contrast_ready.width * 2, contrast_ready.height * 2)
    )

    buffer = BytesIO()
    enlarged.save(buffer, format="PNG")
    return buffer.getvalue()
\end{lstlisting}

The second OCR stage turns raw text into structured JSON
\cite{json2017rfc8259}. The LLM prompt shows that the expected schema
is explicit and validated later through Pydantic
\cite{pydantic2026docs}:

\begin{lstlisting}[language=Python,caption={Nutrition-structuring prompt in \texttt{backend/services/ocr\_service.py}},label={lst:ocr-prompt}]
prompt = f"""
You are given OCR text from a food nutrition label.
Extract only the nutrition information you can confidently identify.

Return only valid JSON in this exact structure:
{{
  "product_name": "string or null",
  "serving_size": "string or null",
  "calories": 0,
  "protein_g": 0,
  "carbs_g": 0,
  "fat_g": 0,
  "sugar_g": 0,
  "sodium_mg": 0,
  "fiber_g": 0
}}
""".strip()
\end{lstlisting}

\subsubsection{Supporting image generation and caching}

Recipe thumbnails are generated only after the textual recipe has
already been accepted. The implementation first derives a deterministic
cache key from normalized ingredient data:

\begin{lstlisting}[language=Python,caption={Thumbnail cache key derivation in \texttt{backend/services/recipe\_image\_service.py}},label={lst:image-cache-key}]
def build_recipe_image_cache_key(
    title: str,
    ingredients: list[str],
    preferences: list[str],
    allergies: list[str],
) -> str:
    normalized_ingredients = sorted(
        _normalize_text(item) for item in ingredients if item.strip()
    )
    digest_source = ",".join(normalized_ingredients)
    return hashlib.sha256(digest_source.encode("utf-8")).hexdigest()[:24]
\end{lstlisting}

The cache lookup and provider delegation are implemented in the same
service:

\begin{lstlisting}[language=Python,caption={Thumbnail reuse and generation logic},label={lst:image-cache-lookup}]
thumbnail_path = _get_thumbnail_path(cache_key)

if thumbnail_path.exists():
    return {
        "image_cache_key": cache_key,
        "image_path": str(thumbnail_path.as_posix()),
        "image_url": _to_public_image_url(cache_key),
        "image_prompt": prompt,
    }

image_bytes = get_recipe_image_provider().generate_png(prompt)
thumbnail_path.write_bytes(image_bytes)
\end{lstlisting}

The concrete provider for image generation is again represented through
an abstract interface and a provider-specific implementation
\cite{openai2026images,openai2026api}:

\begin{lstlisting}[language=Python,caption={Recipe image provider in \texttt{backend/services/recipe\_image\_provider.py}},label={lst:image-provider}]
class RecipeImageProvider(ABC):
    @abstractmethod
    def generate_png(self, prompt: str) -> bytes:
        """Generate a PNG image for the given prompt."""

class OpenAIRecipeImageProvider(RecipeImageProvider):
    def generate_png(self, prompt: str) -> bytes:
        response = self._client.images.generate(
            model=settings.openai_image_model,
            prompt=prompt,
            size="1024x1024",
            quality="low",
            output_format="png",
        )
\end{lstlisting}

This code provides supporting infrastructure for the recipe-generation
feature. Thumbnail generation is not the primary feature outcome, but it
improves the usability of generated results and saved-history views
without changing the main textual workflow.

\subsubsection{Request protection}

Request throttling is implemented in \texttt{backend/rate\_limit.py} as
custom middleware instead of an external gateway dependency. The
middleware keeps an in-memory bucket per client and route:

\begin{lstlisting}[language=Python,caption={Rate-limiting middleware structure},label={lst:rate-limit}]
@dataclass(frozen=True)
class RateLimitRule:
    max_requests: int
    window_seconds: int

class RateLimitMiddleware(BaseHTTPMiddleware):
    async def dispatch(self, request: Request, call_next):
        rule = self._resolve_rule(request)
        client_identifier = self._get_client_identifier(request)
        key = f"{request.method}:{request.url.path}:{client_identifier}"
        allowed, retry_after = self.rate_limiter.check(key, rule)

        if not allowed:
            return JSONResponse(
                status_code=429,
                content={"detail": "Rate limit exceeded. Please try again later."},
                headers={"Retry-After": str(retry_after)},
            )
\end{lstlisting}

The route-specific limits visible earlier in Listing~\ref{lst:main-bootstrap}
show that authentication, recipe generation, and OCR extraction are
protected more strictly than ordinary reads. This is a concrete example
of a cross-cutting concern being realized in middleware rather than
inside every individual route function.

\subsection{Frontend implementation}

\subsubsection{Application shell and routing}

The frontend is organized as a React single-page application whose
top-level router maps directly to the main workflows from the use-case
and navigation diagrams. The routes are declared in
\texttt{frontend/src/App.jsx}, while each major use case is implemented
as its own page component under \texttt{frontend/src/pages/}. This
keeps the frontend aligned with the screen-flow diagram from
Section~3.4 and follows the React and React Router application model
\cite{react2026docs,reactrouter2026docs}.

The page organization is intentionally task-oriented. The Home page
provides entry points into the two primary AI-supported workflows,
Generate and Food Labels. The Login/Account page combines
authentication-related interaction with profile-default management, so
the user does not have to navigate to a separate settings area. The
Saved and History pages then separate two different persistence views:
one for explicitly bookmarked recipes and one for the full generation
history. This mirrors the frontend route structure visible in
\texttt{App.jsx} and keeps each page responsible for one coherent user
task.

\subsubsection{Centralized API access}

HTTP communication \cite{http2022rfc9110} is centralized in
\texttt{frontend/src/lib/api.js}. This module provides a single base URL
resolver and a reusable
\texttt{apiFetch()} wrapper that turns non-success responses into
JavaScript errors with backend-supplied detail messages when available.
This design keeps page components focused on user interaction and local
state instead of repeating low-level request boilerplate. It also gives
the frontend one consistent place to adapt API-base configuration across
development environments.

\subsubsection{Authentication state management}

Shared authentication state is implemented through the React context in
\texttt{frontend/src/context/AuthContext.jsx} \cite{react2026docs}:

\begin{lstlisting}[language=Java,caption={Authentication context in \texttt{frontend/src/context/AuthContext.jsx}},label={lst:auth-context}]
const AuthContext = createContext(null);

export function AuthProvider({ children }) {
  const [token, setToken] = useState(() => localStorage.getItem("token"));
  const [profile, setProfile] = useState(null);
  const [profileLoading, setProfileLoading] = useState(false);

  useEffect(() => {
    if (!token) {
      setProfile(null);
      return;
    }

    const loadProfile = async () => {
      const data = await apiFetch("/me", {
        headers: { Authorization: `Bearer ${token}` },
      });
      setProfile(data);
    };

    loadProfile();
  }, [token]);
}
\end{lstlisting}

This module is the frontend counterpart of the backend
\texttt{get\_current\_user} dependency. The context stores the token,
loads the user profile from \texttt{/me}, and exposes state that other
pages use to decide whether save actions or account-aware defaults are
available.

\subsubsection{Recipe-generation page}

The most interesting frontend behavior is the fallback from manual
fields to stored profile defaults. It is implemented directly in
\texttt{frontend/src/pages/GenerateRecipePage.jsx}:

\begin{lstlisting}[language=Java,caption={Profile-default fallback on the generation page},label={lst:profile-fallback-frontend}]
const manualPreferences = preferences.split(",").map((i) => i.trim()).filter(Boolean);
const manualAllergies = allergies.split(",").map((i) => i.trim()).filter(Boolean);

const payload = {
  ingredients: parsedIngredients,
  preferences: manualPreferences.length > 0 ? manualPreferences : profile?.preferences || [],
  allergies: manualAllergies.length > 0 ? manualAllergies : profile?.allergies || [],
};
\end{lstlisting}

This short block is an important implementation detail because it shows
how profile management and recipe generation are connected across the
frontend and backend. The feature was discussed in the design chapter,
but here it becomes visible as a concrete data-flow decision.

The same page also shows how the generated recipe identifier is used to
persist a bookmarked recipe only after successful generation:

\begin{lstlisting}[language=Java,caption={Saving a generated recipe from the frontend},label={lst:save-recipe-frontend}]
await apiFetch("/recipes", {
  method: "POST",
  headers: {
    "Content-Type": "application/json",
    Authorization: `Bearer ${token}`,
  },
  body: JSON.stringify({ generated_recipe_id: recipe.generated_recipe_id }),
});
\end{lstlisting}

The same page also illustrates the frontend state pattern used
throughout the application: input state is kept locally, asynchronous
actions toggle explicit loading flags, and result or error state is
rendered conditionally. Recipe generation, OCR extraction, profile
updates, and history loading all follow this same approach, which keeps
page behavior easy to reason about without introducing a heavier global
state solution.

\subsubsection{OCR and history-oriented pages}

The Food Labels page combines upload interaction, OCR result rendering,
save behavior, and paginated history loading in a single task-oriented
view. After the user selects an image, the page constructs a
\texttt{FormData} payload for \texttt{/ocr/extract}, renders the raw OCR
text together with the structured nutrition result, and conditionally
enables saving when the user is authenticated. The page also loads
saved OCR history through page-number and page-size query parameters,
which makes the backend pagination contract visible in the frontend.

The Saved Recipes page and Recipe History page follow the same
data-access pattern. Both request paginated backend data, maintain local
\texttt{loading}, \texttt{error}, \texttt{page}, and
\texttt{totalPages} state, and render pagination controls only when more
than one page exists. The difference is semantic rather than structural:
the Saved page filters for explicitly bookmarked recipes, while the
History page displays all generated recipe entries, including unsaved
ones and their warnings. This close structural similarity is useful
because it keeps persistence-related pages consistent for the user while
still preserving their different purposes.

\subsubsection{Structured result presentation}

Structured nutritional output is rendered through the reusable
\texttt{frontend/src/components/NutritionGrid.jsx} component. The
component accepts a nutrition object and an optional field list, which
allows the same grid to display the compact recipe nutrition estimate on
the Generate, Saved, and History pages, and the richer OCR nutrition
schema on the Food Labels page. This is a small but meaningful frontend
abstraction: presentation logic for structured nutrition data is kept in
one component instead of being duplicated across multiple pages.

\subsection{Persistence implementation}

\subsubsection{Relational schema in executable form}

The relational design from Section~\ref{sec:design} is not only diagrammed; it is
materialized through ORM mappings and migration files. The resulting
tables remain small in number but rich enough to capture ownership,
save-state transitions, timestamps, and the denormalized structured
payloads returned by the application. The following DDL-style excerpt
shows the schema shape that the ORM and Alembic revisions realize
\cite{sqlalchemy2026docs,alembic2026docs}:

\begin{lstlisting}[language=SQL,caption={Representative relational DDL for the implemented schema},label={lst:ddl-schema}]
CREATE TABLE users (
    id INTEGER PRIMARY KEY,
    email VARCHAR NOT NULL UNIQUE,
    password VARCHAR NOT NULL,
    preferences TEXT NOT NULL DEFAULT '[]',
    allergies TEXT NOT NULL DEFAULT '[]',
    created_at DATETIME DEFAULT CURRENT_TIMESTAMP
);

CREATE TABLE generated_recipes (
    id INTEGER PRIMARY KEY,
    fingerprint VARCHAR NOT NULL,
    title VARCHAR NOT NULL,
    ingredients TEXT NOT NULL,
    additional_ingredients TEXT NOT NULL DEFAULT '[]',
    preferences TEXT NOT NULL DEFAULT '[]',
    allergies TEXT NOT NULL DEFAULT '[]',
    steps TEXT NOT NULL,
    nutrition TEXT NOT NULL,
    warnings TEXT NOT NULL DEFAULT '[]',
    image_url VARCHAR,
    image_path VARCHAR,
    image_cache_key VARCHAR,
    image_prompt TEXT,
    is_saved BOOLEAN NOT NULL DEFAULT FALSE,
    saved_at DATETIME,
    created_at DATETIME DEFAULT CURRENT_TIMESTAMP,
    user_id INTEGER REFERENCES users(id)
);

CREATE TABLE ocr_extractions (
    id INTEGER PRIMARY KEY,
    source_filename VARCHAR NOT NULL,
    raw_text TEXT NOT NULL,
    structured_nutrition TEXT NOT NULL,
    image_path VARCHAR,
    image_url VARCHAR,
    created_at DATETIME DEFAULT CURRENT_TIMESTAMP,
    user_id INTEGER NOT NULL REFERENCES users(id)
);
\end{lstlisting}

This representation is useful because it makes the key decisions
visible in one place: the ownership model is implemented through foreign
keys, recipe persistence permits a nullable owner because guest
generation exists conceptually, and structured variable-length content
is stored as serialized text rather than decomposed into many auxiliary
tables.

The \texttt{fingerprint} field in \texttt{generated\_recipes} is also an
implementation detail with practical value rather than a decorative
schema element. In \texttt{backend/utils.py}, the application builds
this fingerprint by normalizing the recipe title, ingredients,
preferences, allergies, and steps, serializing the normalized payload,
and hashing it with SHA-256. This makes the identity
of a generated recipe depend on its meaningful content instead of
surface formatting differences such as extra whitespace or input order.
The result is a stable identifier that can be stored together with the
generated row and reused in later comparisons or debugging.

\subsubsection{ORM entities}

The relational model from Section~\ref{sec:design} is implemented in
\texttt{backend/models.py} with SQLAlchemy
\cite{sqlalchemy2026docs}. The \texttt{GeneratedRecipe} entity is the
most important example because it combines ownership, save state,
thumbnail metadata, and several intentionally serialized fields:

\begin{lstlisting}[language=Python,caption={\texttt{GeneratedRecipe} ORM mapping},label={lst:generated-recipe-model}]
class GeneratedRecipe(Base):
    __tablename__ = "generated_recipes"

    id = Column(Integer, primary_key=True, index=True)
    fingerprint = Column(String, nullable=False, index=True)
    title = Column(String, nullable=False)
    ingredients = Column(Text, nullable=False)
    additional_ingredients = Column(Text, nullable=False, default="[]")
    preferences = Column(Text, nullable=False, default="[]")
    allergies = Column(Text, nullable=False, default="[]")
    steps = Column(Text, nullable=False)
    nutrition = Column(Text, nullable=False)
    warnings = Column(Text, nullable=False, default="[]")
    image_url = Column(String, nullable=True)
    image_path = Column(String, nullable=True)
    image_cache_key = Column(String, nullable=True, index=True)
    is_saved = Column(Boolean, nullable=False, default=False, server_default="false")
    saved_at = Column(DateTime(timezone=True), nullable=True)
    created_at = Column(DateTime(timezone=True), server_default=func.now())
    user_id = Column(Integer, ForeignKey("users.id"), nullable=True)

    owner = relationship("User", back_populates="generated_recipes")
\end{lstlisting}

This mapping closes the loop between the ER diagram, the relational
schema, and the implementation. The JSON-serialized text fields are a
deliberate denormalization tradeoff already justified in
Section~3.3.4: they simplify storage of variable-length ingredient,
warning, and step lists without introducing multiple auxiliary tables
for a thesis-scale application.

\subsubsection{Schema migration history}

Schema evolution is managed through Alembic \cite{alembic2026docs}. The migration
\texttt{a3f9c2d1e804\_ocr\_image\_fields\_and\_pagination\_indexes.py}
is a good example because it adds both new columns and composite
indexes:

\begin{lstlisting}[language=Python,caption={Alembic revision adding image fields and pagination indexes},label={lst:alembic-indexes}]
def upgrade() -> None:
    op.add_column("ocr_extractions", sa.Column("image_path", sa.String(), nullable=True))
    op.add_column("ocr_extractions", sa.Column("image_url", sa.String(), nullable=True))

    with op.batch_alter_table("generated_recipes") as batch_op:
        batch_op.create_index(
            "ix_generated_recipes_user_saved_at",
            ["user_id", "is_saved", "saved_at"],
        )
        batch_op.create_index(
            "ix_generated_recipes_user_created_at",
            ["user_id", "created_at"],
        )
\end{lstlisting}

This revision demonstrates that the relational design did not remain
static after the ERD was drawn. The database schema evolved together
with the implementation needs of OCR history and paginated history
views.

The same evolution also shows where the implementation became more
precise than the first conceptual sketch. The design stabilized around
three adjustments in particular. First, recipe history and explicit
saved recipes were kept in the same table instead of being split into
two separate persistence models, because save-state is an attribute of a
generated result rather than a different domain object. Second, OCR
history gained image-path and image-URL fields only once the frontend
history view made image recall part of the user-visible workflow.
Third, pagination support led directly to composite indexes on
\texttt{generated\_recipes} and \texttt{ocr\_extractions}, so the final
schema reflects not only domain entities but also the access patterns
that emerged during implementation.

\subsubsection{Query and indexing implementation}

The composite indexes above are used by the saved-recipes and history
queries in \texttt{backend/routes/recipes.py}. The saved-recipes query
is shown below:

\begin{lstlisting}[language=Python,caption={Saved-recipes query using the ownership and save-state fields},label={lst:saved-query}]
base_query = db.query(GeneratedRecipe).filter(
    GeneratedRecipe.user_id == current_user.id,
    GeneratedRecipe.is_saved.is_(True),
)
rows = (
    base_query
    .order_by(GeneratedRecipe.saved_at.desc(), GeneratedRecipe.created_at.desc())
    .offset((page - 1) * page_size)
    .limit(page_size)
    .all()
)
\end{lstlisting}

This query is not only persistence code but also evidence that the
index strategy described in Section~\ref{sec:design} is connected to an actual access
pattern in the system.

The same query can be expressed at SQL level as follows:

\begin{lstlisting}[language=SQL,caption={SQL form of the saved-recipes listing query},label={lst:saved-query-sql}]
SELECT
    id,
    title,
    ingredients,
    additional_ingredients,
    preferences,
    allergies,
    steps,
    nutrition,
    image_url,
    image_cache_key,
    saved_at,
    created_at
FROM generated_recipes
WHERE user_id = :current_user_id
  AND is_saved = TRUE
ORDER BY saved_at DESC, created_at DESC
LIMIT :page_size OFFSET :offset;
\end{lstlisting}

This query uses the same ownership and save-state conditions that
appear in the ORM code and benefits directly from the composite index on
\texttt{(user\_id, is\_saved, saved\_at)}. The ordering also reflects a
domain choice: once a recipe becomes part of the saved collection, the
most relevant timestamp is the save action rather than the original
generation time.

Although the user-facing application mainly executes ownership-scoped
single-table reads, the same relational model also supports more
analytical queries when needed. For example, the following aggregation
query can summarize per-user recipe activity by combining a join,
conditional aggregation, and descending ranking:

\begin{lstlisting}[language=SQL,caption={Analytical SQL over the implemented recipe-history schema},label={lst:activity-aggregation-sql}]
SELECT
    u.email,
    COUNT(gr.id) AS generated_recipe_count,
    SUM(CASE WHEN gr.is_saved = TRUE THEN 1 ELSE 0 END) AS saved_recipe_count,
    MAX(gr.created_at) AS latest_generation_at
FROM users AS u
LEFT JOIN generated_recipes AS gr
    ON gr.user_id = u.id
GROUP BY u.id, u.email
ORDER BY generated_recipe_count DESC, latest_generation_at DESC;
\end{lstlisting}

This query is not part of the current browser workflow, but it is still
important documentation evidence because it shows that the schema
supports not only transactional reads and writes but also higher-level
reporting over the same ownership model.

The retrieval corpus is indexed offline through
\texttt{backend/scripts/index\_recipes.py}:

\begin{lstlisting}[language=Python,caption={Offline recipe-corpus indexing for ChromaDB},label={lst:index-script}]
def index_recipe_corpus() -> None:
    recipes = load_recipe_corpus(corpus_path)
    collection = get_collection()
    existing = collection.get(include=[])
    existing_ids = existing.get("ids", [])
    if existing_ids:
        collection.delete(ids=existing_ids)

    for start in range(0, len(recipes), batch_size):
        batch = recipes[start : start + batch_size]
        documents = [build_recipe_document(recipe) for recipe in batch]
        embeddings = create_embeddings(documents)
        collection.add(
            ids=[str(recipe["recipe_id"]) for recipe in batch],
            documents=documents,
            embeddings=embeddings,
            metadatas=[build_recipe_metadata(recipe) for recipe in batch],
        )
\end{lstlisting}

The indexed recipe corpus is based on a cleaned subset of the Food.com
dataset obtained from Kaggle and stored locally in
\texttt{backend/data/recipes\_subset.json} \cite{foodcom2019kaggle}.
This indexing script shows that the retrieval layer is not abstract or
purely conceptual. The corpus preparation and embedding workflow
\cite{openai2026embeddings} are
implemented as executable code that can rebuild the persistent Chroma
store \cite{chromadb2026docs} used during generation.

However, the indexing script is only the second half of the retrieval
pipeline. Before any embeddings are created, the raw Food.com export is
cleaned by \texttt{backend/scripts/prepare\_recipes.py}. This script
parses the CSV source, decodes list-like fields, normalizes ingredient
text, removes duplicates, filters invalid recipes, and builds a short
instruction summary for each retained item. The resulting
\texttt{recipes\_subset.json} file is therefore not a raw dataset dump
but a project-specific intermediate corpus shaped specifically for later
retrieval. This preprocessing step is important implementation evidence
because the quality of retrieval depends not only on the embedding model
but also on how consistently the local corpus is cleaned and reduced
before it enters Chroma.

The vector-store layer itself is implemented in
\texttt{backend/services/vector\_store\_service.py}. That module creates
the persistent Chroma collection, converts each recipe into a compact
document containing the title, ingredients, tags, and instruction
summary, and stores the richer structured fields as metadata alongside
the embedding \cite{chromadb2026docs}. This means that the retrieval
representation used by the project is deliberately asymmetric: the
document text is optimized for semantic matching, while metadata keeps
the structured recipe fields available for later filtering, ranking, and
prompt construction.

\subsection{Known implementation limitations}

Several implementation limitations remain and are important for honest
evaluation.

\subsubsection{Known bugs and unresolved issues}

The first unresolved issue is that the quality of recipe generation and
nutrition-label extraction depends partly on upstream AI-provider
behavior. Even when the local validation and control flow operate
correctly, the application may still receive weak OCR text, malformed
model output, or structurally poor generated content from the external
services it relies on. In such cases, the backend may reject the
response, return a lower-quality result, or surface warnings instead of
producing a fully satisfactory output. This is not a control-flow bug in
the local application logic, but it remains an unresolved reliability
issue because output quality cannot be guaranteed purely by local code.

A second unresolved issue concerns OCR upload persistence. In
\texttt{backend/services/ocr\_service.py}, the system attempts to store
the uploaded image before OCR extraction begins, but a failure during
that media-write step does not stop the extraction workflow. As a
result, OCR processing may still succeed and return structured nutrition
data even when the uploaded image itself was not persisted correctly.
This behavior keeps the core extraction feature available in the
presence of non-critical storage errors, but it also means that a saved
history entry may lack a usable stored-image reference.

A third practical weakness is the limited recovery strategy around
AI-generated recipe validity. The recipe workflow validates generated
content against restriction rules and performs one corrective retry when
the first result violates allergy or dietary constraints. However, if
both attempts fail, the request is rejected. This keeps the safety model
strict, but it also means that transient or provider-specific output
instability can still cause complete request failure even when the input
itself is valid.

\subsubsection{Technical debt and maintainability concerns}

The default relational database is SQLite \cite{sqlite2026docs}, as configured in
\texttt{backend/config.py}. This is appropriate for a local thesis-scale
deployment, but it is not designed for larger multi-user or
high-concurrency scenarios. The same applies to the custom rate limiter
in \texttt{backend/rate\_limit.py}, which stores counters in process
memory. Its limits are therefore reset when the backend restarts and are
not shared across multiple application instances.

Authentication state on the frontend is currently stored in browser
\texttt{localStorage}, as implemented in
\texttt{frontend/src/context/AuthContext.jsx}. This keeps the frontend
logic simple, but it is a weaker session-management approach than using
refresh tokens, httpOnly cookies, or server-managed session renewal.
Similarly, account management remains intentionally minimal: the current
implementation supports signup, login, logout, and profile-default
updates, but it does not include password reset, email verification, or
role-based authorization features.

The persistence model also contains an intentional denormalization
tradeoff. Several structured values, such as recipe ingredients, steps,
warnings, and OCR nutrition payloads, are stored as serialized text
rather than being decomposed into separate relational child tables. This
keeps the schema compact and matches the JSON-shaped API contracts
\cite{json2017rfc8259} well,
but it reduces query flexibility and would become harder to maintain if
the system evolved toward more complex reporting or cross-record
analytics.

Media handling is likewise optimized for simplicity rather than for a
strong deployment model. OCR upload images and generated recipe
thumbnails are stored on the local filesystem under the backend
workspace. This is sufficient for local development and demonstration,
but it would need to be replaced by a more robust object-storage
strategy in a larger deployment.

The thumbnail cache also reflects a pragmatic compromise. In
\texttt{backend/services/recipe\_image\_service.py}, the cache key is
derived mainly from normalized ingredient input. This reduces repeated
image-generation calls and API cost, but it can also cause requests with
the same base ingredient set to reuse the same thumbnail even when
preference or allergy context differs.

The application also remains dependent on external AI services for its
main intelligent features. Recipe generation, embedding creation
\cite{openai2026embeddings}, OCR text structuring, and recipe-image
generation all require upstream API access and therefore do not
currently support a fully offline execution mode. This is acceptable
for the current thesis-scale deployment model, but it defines a clear
operational boundary for the implementation.

\subsubsection{Planned but unimplemented improvements}

Several improvements were deliberately left outside the thesis scope.
The recipe workflow does not yet support editing and re-saving generated
recipes after the initial result is stored. The system does not yet
provide editing or deletion for saved recipes or saved OCR history
entries. The OCR subsystem exposes a provider abstraction, but only the
Google Vision-based implementation \cite{google2026vision} is currently
available in production form. Likewise, recipe image generation is
represented through a provider interface, but only the OpenAI-backed
path \cite{openai2026images,openai2026api} is implemented.

The current application also lacks more advanced administrative and
analytical functionality. For example, there is no dedicated dashboard
for aggregate usage statistics, no administrator view over user activity
patterns, and no richer moderation or review workflow for generated
content. These omissions do not prevent the core thesis goals from being
demonstrated, but they remain natural extension points.

\subsubsection{Future development directions}

Future work should first strengthen operational robustness. The most
important infrastructure-oriented improvements would be replacing the
in-memory rate limiter with a shared backing store, migrating from
SQLite to a stronger multi-user database when needed, and making media
persistence more explicit and failure-aware.

On the AI side, future development should broaden provider support,
improve evaluation coverage, and strengthen output control. Useful next
steps would include adding alternative OCR backends, expanding the
recipe and OCR evaluation datasets, refining retrieval and prompt
construction, and introducing more flexible post-generation revision
controls for users.

From a product perspective, the most valuable user-facing improvements
would be edit/delete support for stored content, stronger account
management, and clearer lifecycle handling for generated artifacts and
saved OCR media. These directions follow directly from the current
implementation and would extend the project without requiring a
fundamental redesign.

\subsection{Section Summary}

The design section presented the architecture and design rationale of the
project; this implementation section showed their realization in code.
The FastAPI bootstrap, route modules, service classes, and React
context/page structure demonstrate that the UML and architectural
diagrams reflect the actual repository. The persistence layer, Alembic
revisions, and Chroma indexing script likewise show that the database
and retrieval designs were implemented as executable program components
rather than remaining design-only artifacts. This concrete
implementation basis is used in the testing section to evaluate
functional behavior, robustness, and output quality.
